\documentclass[12pt]{article}
\newcommand{\DocShort}{Método geométrico}
\input{preamble_population.tex}
\begin{document}
\doccover{Método del incremento geométrico}{2026-05}
\handoutintro{Proyectar la población futura suponiendo que la tasa de crecimiento relativa por período censal permanece constante.}{Ciudades nuevas o en expansión, especialmente cuando el crecimiento relativo todavía es alto.}

\section{Datos ejemplo}
Se usará una serie con períodos censales de $10$ años; por tanto, las tasas calculadas en el ejemplo son tasas por década.

\begin{center}
\begin{tabular}{lcccccc}
\toprule
Año & 1970 & 1980 & 1990 & 2000 & 2010 & 2020 \\
\midrule
Población & 45,200 & 53,400 & 63,300 & 75,100 & 88,700 & 104,900 \\
\bottomrule
\end{tabular}
\end{center}

\section{Ecuaciones mínimas}
La tasa relativa del período censal $i$ es:
\begin{equation}
r_i=\frac{P_i-P_{i-1}}{P_{i-1}}
\end{equation}
donde $P_{i-1}$ es la población al inicio del período y $P_i$ es la población al final del período.

La tasa geométrica media por período censal es:
\begin{equation}
I_G=\left(r_1r_2\cdots r_N\right)^{1/N}
\end{equation}
donde $I_G$ es la tasa geométrica media del conjunto y $N$ es el número de períodos.

La población proyectada después de $n$ períodos censales es:
\begin{equation}
P_n=P_{\text{actual}}(1+I_G)^n
\end{equation}
donde $P_{\text{actual}}$ es la población del último censo y $n$ es el número de períodos futuros.

\weightednote{Si los intervalos censales no son regulares, conviene convertir cada tasa a base anual y usar un promedio ponderado por duración del período. Una opción práctica es definir la tasa anual efectiva
\[
g_i=\left(\frac{P_i}{P_{i-1}}\right)^{1/\Delta t_i}-1
\]
y luego promediarla con pesos $\Delta t_i$ en años.}

\section{Procedimiento}
\begin{enumerate}
\item Calcular el incremento de cada período censal.
\item Convertir cada incremento en tasa relativa respecto a la población de inicio del período.
\item Obtener la media geométrica de esas tasas por período.
\item Aplicar la ecuación de crecimiento compuesto.
\end{enumerate}

\section{Ejemplo resuelto}
\begin{center}
\begin{tabular}{crrr}
\toprule
Año & Población & Incremento del período & Tasa relativa por década \\
\midrule
1970 & 45,200 & -- & -- \\
1980 & 53,400 & 8,200 & 0.1814 \\
1990 & 63,300 & 9,900 & 0.1854 \\
2000 & 75,100 & 11,800 & 0.1864 \\
2010 & 88,700 & 13,600 & 0.1811 \\
2020 & 104,900 & 16,200 & 0.1826 \\
\bottomrule
\end{tabular}
\end{center}

\[
I_G=(0.1814\times0.1854\times0.1864\times0.1811\times0.1826)^{1/5}=0.1834
\]

\[
P_{2030}=\num{104900}(1+0.1834)^1=\num{124136}
\]
\[
P_{2040}=\num{104900}(1+0.1834)^2=\num{146900}
\]
\[
P_{2050}=\num{104900}(1+0.1834)^3=\num{173839}
\]

\resultbox{La población proyectada es 124,136 habitantes en 2030, 146,900 en 2040 y 173,839 en 2050.}

\section{Teoría breve}
Este método representa crecimiento compuesto. En vez de mantener constante el aumento absoluto, mantiene constante el aumento relativo por período censal.

\section{Observaciones de uso}
\begin{itemize}
\item Produce valores más altos que el método aritmético.
\item Puede sobreestimar si la ciudad ya se acerca a su saturación.
\item Debe aclararse si la tasa está expresada por año, por década o por otro período.
\end{itemize}
\end{document}
